%
% architecture.tex
%
% Beadle: Architecture & Z Specification
% Package structure, trust model, credential resolution, message lifecycle,
% and security boundaries for beadle-email and beadle-daemon.
%
% Type-check: fuzz -t architecture.tex
%

\documentclass[a4paper,10pt,fleqn]{article}
\usepackage[margin=1in]{geometry}
\usepackage{fuzz}
\usepackage[colorlinks=true,linkcolor=blue,citecolor=blue,urlcolor=blue]{hyperref}
\usepackage{booktabs}
\usepackage{longtable}

\begin{document}

\title{Beadle: Architecture \& Z Specification}
\author{Punt Labs}
\date{March 2026}
\hypersetup{
  pdftitle={Beadle: Architecture and Z Specification},
  pdfauthor={Punt Labs},
  pdfsubject={Architecture Document},
  pdfcreator={fuzz}
}
\maketitle

\tableofcontents
\newpage

%%----------------------------------------------------------------------
\section{Introduction}
%%----------------------------------------------------------------------

Beadle is an autonomous agent daemon with cryptographic owner control.
Command \emph{execution} requires a GPG-signed command file authorized
by the owner (\S~Command-File Signature Verification); the daemon has
no independent decision-making about \emph{what to run}. The earlier
steps that decide whether to create a pipeline at all --- sender
permission and transport-trust checks --- carry no such signature (\S~Daemon
Architecture, Design Invariant~1). ``Owner'' here names the single GPG
key \texttt{daemon.json} designates as the recipe authorizer (via
\texttt{owner\_handle} or \texttt{owner\_gpg\_key\_id}), not a general
position of authority over the system; a rename to ``authorizer'' is
proposed under bead beadle-4ul but not yet done.

\texttt{beadle-email} shipped first: an MCP server and CLI that provide
email communication via Proton Bridge with a four-level PGP trust
model. \texttt{beadle-daemon} is a second binary, sharing
\texttt{beadle-email}'s lower package layers, that acts on that mail
unattended (\S~System Overview, \S~Daemon Architecture). This document
describes both binaries: their package structure, data flows, security
boundaries, and key state machines formalised in Z.

%%----------------------------------------------------------------------
\section{System Overview}
%%----------------------------------------------------------------------

Beadle ships as two independent binaries sharing package layers.
\texttt{beadle-email} exposes an MCP server and CLI for message
send/receive/triage; \texttt{beadle-daemon} watches a mailbox
unattended and turns authorized instructions into mission pipelines.
Neither imports the other at the command layer (\S~Dependency Graph).

\subsection{beadle-email}

\texttt{beadle-email} runs as an MCP server, by default over stdio;
\texttt{serve --transport ws} instead starts a WebSocket listener
(\S~Security Boundaries).  It connects to Proton Bridge via IMAP on
localhost by default --- a config default, not an enforced boundary
(\S~Security Boundaries) --- sends mail via Proton Bridge SMTP
(primary) or Resend API (fallback), and verifies PGP signatures using
an isolated GPG keyring.

The binary groups its subcommands by concern: server administration
(\texttt{serve}, \texttt{doctor}, \texttt{status}, \texttt{version},
\texttt{health}), message operations (\texttt{list}, \texttt{search},
\texttt{read}, \texttt{send}, \texttt{reply}, \texttt{move},
\texttt{mark}, \texttt{folders}), contact management
(\texttt{contact list|add|remove|find}), identity
(\texttt{identity show|set}), and enablement (\texttt{enable}/
\texttt{disable}, \texttt{install}/\texttt{uninstall}).  Run
\texttt{beadle-email --help} for the current, authoritative list ---
it changes independently of this document.

\subsection{beadle-daemon}

\texttt{beadle-daemon} has no MCP surface at all --- it does not
import \texttt{internal/mcp} (\S~Dependency Graph).  Its single
subcommand, \texttt{run}, polls the mailbox, checks sender permission
and transport trust (\S~Daemon Architecture), and on a pass executes a
signed-command pipeline.  It is implemented but not released:
\texttt{make build-daemon} builds it locally, but no CI workflow and
no \texttt{make dist} target produces a distributable binary (bead
beadle-5ma).

%%----------------------------------------------------------------------
\section{Package Map}
%%----------------------------------------------------------------------

Beadle's code lives under \texttt{cmd/} (binary entry points) and
\texttt{internal/} (everything else --- nothing under
\texttt{internal/} is exported outside the module). Three further
\texttt{internal/} directories --- \texttt{daemontest},
\texttt{testenv}, \texttt{testserver} --- are test-only fixtures for
the suite and are excluded from the table below, which covers the
remaining production packages.

\begin{center}
\begin{tabular}{@{}ll@{}}
\toprule
\textbf{Package} & \textbf{Responsibility} \\
\midrule
\texttt{cmd/beadle-daemon/}  & Daemon entry point: the mail-triggered mission pipeline runner (\texttt{run}) \\
\texttt{cmd/beadle-email/}   & CLI/MCP entry point: subcommand dispatch (\S~System Overview) \\
\texttt{internal/channel/}   & Channel interface: \texttt{Message}, \texttt{TrustLevel}, shared types \\
\texttt{internal/claudemd/}  & Adds/removes beadle's \texttt{@}-import line in a repo's \texttt{CLAUDE.md} \\
\texttt{internal/contacts/}  & Address book storage and lookup \\
\texttt{internal/daemon/}    & Mail-triggered mission pipeline: planner, executor, runners, command model, persistence \\
\texttt{internal/email/}     & IMAP client, MIME parser, trust classifier, SMTP/Resend senders \\
\texttt{internal/enable/}    & Turns beadle's \texttt{CLAUDE.md} guidance composition on/off in a repo \\
\texttt{internal/identity/}  & Resolves which identity beadle operates as \\
\texttt{internal/mcp/}       & MCP tool definitions and handlers (\S~MCP Tool Surface); stdio and WebSocket transports \\
\texttt{internal/paths/}     & Single root directory for all beadle data \\
\texttt{internal/pgp/}       & GPG signature verification (isolated GNUPGHOME) and signing (system keyring) \\
\texttt{internal/secret/}    & Credential resolution: OS keychain $\to$ file $\to$ env var \\
\texttt{internal/session/}   & Reads the ethos session roster to enumerate participants \\
\bottomrule
\end{tabular}
\end{center}

\subsection{Dependency Graph}

Dependencies flow strictly downward within each binary; the edges
below are the internal import edges reported by \texttt{go list}, not
a hand-traced approximation.  \texttt{internal/mcp} and
\texttt{internal/daemon} do not import each other --- the two binaries
share a lower layer (\texttt{channel}, \texttt{contacts},
\texttt{email}, \texttt{identity}, \texttt{paths}, \texttt{pgp},
\texttt{secret}) and are independent at the command layer, matching
the split in \S~System Overview.

\begin{verbatim}
beadle-email command layer:
  cmd/beadle-email -> mcp, contacts, email, enable, identity,
                       paths, pgp, secret
  mcp              -> channel, contacts, email, enable, identity,
                       paths, pgp, session
  enable           -> claudemd

beadle-daemon command layer:
  cmd/beadle-daemon -> daemon, email, identity, paths, secret
  daemon            -> channel, contacts, email, identity, paths, pgp

shared lower layer (imported by both trees above):
  email    -> channel, contacts, identity, paths, pgp, secret
  contacts -> paths
  secret   -> paths
  channel, claudemd, identity, paths, pgp, session -> (no internal deps)
\end{verbatim}

%%----------------------------------------------------------------------
\section{Basic Types}
%%----------------------------------------------------------------------

\begin{zed}
  [MSGID, ADDRESS, BODY, HEADER, RAWBYTES, KEYID, SIGNER, CREDNAME]
\end{zed}

$MSGID$ is an IMAP UID (numeric string).  $ADDRESS$ is an RFC~5322
email address.  $BODY$ is the plain-text message body.  $HEADER$ is
a header name-value pair.  $RAWBYTES$ is the raw RFC~822 message.
$KEYID$ is a GPG key fingerprint.  $SIGNER$ is a GPG user identity.
$CREDNAME$ is a named credential (e.g.\ \texttt{imap-password}).

%%----------------------------------------------------------------------
\section{Free Types}
%%----------------------------------------------------------------------

\begin{zed}
  ZBOOL ::= ztrue | zfalse
\end{zed}

\begin{zed}
  TrustLevel ::= trusted | verified | untrusted | unverified
\end{zed}

Four trust levels classify inbound messages:

\begin{itemize}
  \item $trusted$ --- Proton-to-Proton, end-to-end encrypted by Proton
    infrastructure.
  \item $verified$ --- external sender with a valid PGP signature.
  \item $untrusted$ --- external sender with an invalid PGP signature.
  \item $unverified$ --- external sender with no signature present.
\end{itemize}

\begin{zed}
  SendMethod ::= smProtonSMTP | smResend
\end{zed}

Outbound messages use Proton Bridge SMTP as the primary transport,
falling back to Resend API when Bridge is unavailable.

\begin{zed}
  CredBackend ::= cbKeychain | cbFile | cbEnvVar
\end{zed}

Credentials are resolved through a priority chain of three backends.

\begin{zed}
  IMAPState ::= imapDisconnected | imapConnected | imapSelected
\end{zed}

The IMAP client progresses through connection, authentication, and
folder selection before fetching messages.

%%----------------------------------------------------------------------
\section{Trust Model}
%%----------------------------------------------------------------------

This section formalises the trust classification state machine.
Trust classification is a pure function of message headers and
PGP verification outcome.

\begin{schema}{TrustInput}
  hasProtonE2E : ZBOOL \\
  hasProtonInternal : ZBOOL \\
  hasPGPSignature : ZBOOL \\
  pgpVerified : ZBOOL
\end{schema}

$hasProtonE2E$ is true when the header
\texttt{X-Pm-Content-Encryption} contains \texttt{end-to-end}.
$hasProtonInternal$ is true when \texttt{X-Pm-Origin} is
\texttt{internal}.  $hasPGPSignature$ is true when the Content-Type
is \texttt{multipart/signed} or raw bytes contain
\texttt{application/pgp-signature}.  $pgpVerified$ is the result of
\texttt{gpg --verify} (only meaningful when $hasPGPSignature = ztrue$).

\begin{schema}{ClassifyTrust}
  TrustInput \\
  level! : TrustLevel
  \where
  (hasProtonE2E = ztrue \land hasProtonInternal = ztrue
    \implies level! = trusted) \\
  \land \\
  (hasPGPSignature = ztrue \land pgpVerified = ztrue
    \land \lnot (hasProtonE2E = ztrue \land hasProtonInternal = ztrue)
    \implies level! = verified) \\
  \land \\
  (hasPGPSignature = ztrue \land pgpVerified = zfalse
    \land \lnot (hasProtonE2E = ztrue \land hasProtonInternal = ztrue)
    \implies level! = untrusted) \\
  \land \\
  (hasPGPSignature = zfalse
    \land \lnot (hasProtonE2E = ztrue \land hasProtonInternal = ztrue)
    \implies level! = unverified)
\end{schema}

The classification is total: every combination of inputs maps to
exactly one trust level.  Proton E2E takes priority over PGP
signatures (a Proton-internal message with a PGP signature is
still $trusted$, not $verified$).

\subsection{Two-Phase Trust Resolution}

Trust classification happens in two phases because PGP verification
requires fetching the raw message bytes and running \texttt{gpg}:

\begin{enumerate}
  \item \textbf{Header scan} --- \texttt{ClassifyTrust} runs with
    $pgpVerified = zfalse$.  Messages with signatures return
    $unverified$ with $hasPGPSignature = ztrue$, signalling the
    caller to proceed to phase~2.
  \item \textbf{PGP verification} --- the caller fetches raw bytes,
    runs \texttt{gpg --verify} in an isolated GNUPGHOME, and
    re-classifies with the actual $pgpVerified$ value.
\end{enumerate}

%%----------------------------------------------------------------------
\section{Credential Resolution}
%%----------------------------------------------------------------------

\begin{schema}{CredentialStore}
  keychainAvail : ZBOOL \\
  keychainValues : CREDNAME \pfun BODY \\
  fileValues : CREDNAME \pfun BODY \\
  envValues : CREDNAME \pfun BODY
\end{schema}

Credentials are resolved by name.  The keychain backend maps to
OS-specific stores.  On Darwin, the macOS Keychain via the
\texttt{security} CLI, keyed by (service=\texttt{beadle},
account=\texttt{<name>}).  On Linux, a two-stage lookup: first
\texttt{pass} (\texttt{pass show beadle/<name>}), which keeps
secrets GPG-encrypted with the user's own key and matches the
Proton Bridge vault backend; then \texttt{secret-tool} (libsecret
/ GNOME Keyring), keyed by the same (service, account) pair as
Darwin for portability.  If both Linux binaries are present,
\texttt{pass} wins.  File credentials live in
\texttt{\~{}/.punt-labs/beadle/secrets/<name>} with mode 600.
Environment variables follow the pattern \texttt{BEADLE\_<NAME>}.

\begin{schema}{ResolveCredential}
  \Xi CredentialStore \\
  name? : CREDNAME \\
  value! : BODY
  \where
  (keychainAvail = ztrue \land name? \in \dom keychainValues
    \implies value! = keychainValues(name?)) \\
  \land \\
  ((\lnot (keychainAvail = ztrue \land name? \in \dom keychainValues))
    \land name? \in \dom fileValues
    \implies value! = fileValues(name?)) \\
  \land \\
  ((\lnot (keychainAvail = ztrue \land name? \in \dom keychainValues))
    \land name? \notin \dom fileValues
    \land name? \in \dom envValues
    \implies value! = envValues(name?))
\end{schema}

The resolution is a priority chain: keychain beats file, file beats
environment variable.  File credentials are rejected if group or
world readable (permissions check enforces mode 600).

The credentials used by \texttt{beadle-email}, resolved by name:

\begin{center}
\begin{tabular}{@{}lll@{}}
\toprule
\textbf{Name} & \textbf{Env Var} & \textbf{Purpose} \\
\midrule
\texttt{imap-password}   & \texttt{BEADLE\_IMAP\_PASSWORD}   & Proton Bridge login \\
\texttt{smtp-password}   & \texttt{BEADLE\_SMTP\_PASSWORD}   & SMTP login; falls back to \texttt{imap-password} when unset \\
\texttt{resend-api-key}  & \texttt{BEADLE\_RESEND\_API\_KEY}  & Resend API authentication \\
\texttt{gpg-passphrase}  & \texttt{BEADLE\_GPG\_PASSPHRASE}  & PGP signing key passphrase \\
\bottomrule
\end{tabular}
\end{center}

A secret file that exists but is unreadable or malformed (permissions,
I/O error) is a hard error --- resolution does not fall through to the
environment variable in that case, only when the file is genuinely
absent.

%%----------------------------------------------------------------------
\section{Message Receive Path}
%%----------------------------------------------------------------------

The receive path is the data flow from IMAP mailbox to MCP tool
response.

\begin{schema}{IMAPConnection}
  state : IMAPState \\
  host : ADDRESS \\
  user : ADDRESS
\end{schema}

\begin{schema}{Connect}
  \Delta IMAPConnection
  \where
  state = imapDisconnected \\
  state' = imapConnected
\end{schema}

Connection uses STARTTLS.  The TLS certificate check is skipped
automatically for loopback addresses (Proton Bridge uses self-signed
certificates on localhost); \texttt{tls\_skip\_verify} in config
extends the same skip to any host (Design Invariant~10).

\begin{schema}{SelectFolder}
  \Delta IMAPConnection \\
  folder? : CREDNAME
  \where
  state = imapConnected \lor state = imapSelected \\
  state' = imapSelected
\end{schema}

\subsection{Fetch and Parse}

Once a folder is selected, messages are fetched by UID.  The raw
RFC~822 bytes are parsed through a pipeline:

\begin{enumerate}
  \item \textbf{PGP decrypt} (conditional) --- if the message is
    PGP-encrypted, inbound decryption (DES-024, SETTLED) runs
    \emph{before} MIME parsing, using the system keyring (a
    deliberate departure from the isolated-keyring rule for
    verification, since decryption needs the owner's own private
    key --- Design Invariant~3).  This changes which bytes MIME
    parse sees.
  \item \textbf{MIME parse} --- \texttt{ParseMIME} extracts the
    plain-text body, attachments, and selected headers.  Multipart
    messages are walked using the \texttt{go-message} library, and
    \texttt{walkParts} recurses into a nested \texttt{multipart} part
    before reading its body (DES-025, SETTLED) --- GPG Mail nests
    \texttt{multipart/signed} $\to$ \texttt{multipart/mixed} $\to$
    \texttt{text/plain}, and a flat, non-recursive walk returned
    ``(no text body)'' for exactly that shape.  PGP signature and key
    parts are skipped (handled separately).
  \item \textbf{Header trust} --- \texttt{ClassifyTrust} runs on
    headers to determine the initial trust level.
  \item \textbf{PGP verify} (conditional) --- if the message has a
    PGP signature ($hasPGPSignature = ztrue$), raw bytes are
    passed to \texttt{pgp.Verify} for full verification; the
    verify-and-promote step that turns this into a final trust level
    lives in the MCP handlers, not in \texttt{internal/email} --- see
    \S~Two-Phase Trust Resolution.
  \item \textbf{Response assembly} --- the parsed message, trust
    level, and attachments are serialised as JSON in the MCP tool
    response.
\end{enumerate}

%%----------------------------------------------------------------------
\section{Message Send Path}
%%----------------------------------------------------------------------

\begin{schema}{SendState}
  smtpAvailable : ZBOOL \\
  method! : SendMethod
  \where
  smtpAvailable = ztrue \implies method! = smProtonSMTP \\
  smtpAvailable = zfalse \implies method! = smResend
\end{schema}

The send path attempts Proton Bridge SMTP first.  If the Bridge is
unreachable (a TCP dial to the SMTP port times out in 2~seconds),
the Resend API is used as fallback.

\subsection{Proton Bridge SMTP}

SMTP uses implicit TLS on port 465, or STARTTLS otherwise, with the
same self-signed certificate exception as IMAP for loopback addresses
--- extended to any host when \texttt{tls\_skip\_verify} is set
(Design Invariant~10).  Port 465 implicit TLS is also the branch the
Fastmail signing transport uses (\S~PGP Signing).  Authentication uses
\texttt{PLAIN} with \texttt{smtp-password} if configured, else the
same credentials as IMAP (\S~Credential Resolution).  The raw RFC~822
message is composed by \texttt{ComposeRaw}, which validates header
fields against CR/LF injection.

\subsection{Resend API}

The Resend fallback sends via HTTPS POST to
\texttt{api.resend.com/emails}.  It supports both plain text and
HTML bodies.  The API key is resolved via the credential store.

\subsection{PGP Signing}

PGP signing of outbound mail (DES-022, SETTLED) requires a transport
that preserves raw MIME (\texttt{multipart/signed} envelopes).
Neither Proton Bridge (strips MIME) nor Resend (no raw MIME API;
Resend's own backend is Amazon SES, which strips it the same way)
supports this.  The signing transport that does is \textbf{Fastmail
SMTP} on port 465 (implicit TLS, \S~Proton Bridge SMTP) --- selected
by pointing \texttt{smtp\_host}/\texttt{smtp\_port} at Fastmail (see
\texttt{docs/TESTING.md}'s Fastmail Test Config).

\texttt{TrySendChain} does not itself validate that the transport it
used preserved the signed MIME envelope --- it reports success on any
clean SMTP transaction.  A signed message sent over the default
(Proton Bridge) configuration reports ``sent'' while carrying a
signature Bridge has already stripped; only a transport actually known
to preserve \texttt{multipart/signed}, such as Fastmail, makes the
signature meaningful on arrival.

%%----------------------------------------------------------------------
\section{PGP Subsystem}
%%----------------------------------------------------------------------

The PGP subsystem handles signature verification and signing via
the system \texttt{gpg} binary.  It never uses a Go PGP library ---
\texttt{gpg} handles keyring management, trust models, and
signature formats.

\begin{schema}{GPGEnvironment}
  systemKeys : \power KEYID \\
  importedKeys : \power KEYID \\
  isolated : ZBOOL
  \where
  isolated = ztrue \\
  importedKeys \subseteq systemKeys
\end{schema}

Every verification runs in an isolated GNUPGHOME under
\texttt{/tmp/bg-*}.  The short path prefix avoids the 108-byte
Unix socket path limit for \texttt{gpg-agent}.

\subsection{Verify Flow}

\begin{enumerate}
  \item Extract signed data, detached signature, and optional public
    key from the \texttt{multipart/signed} envelope (RFC~3156).
  \item Create an isolated GNUPGHOME directory.
  \item If the message includes a public key attachment
    (\texttt{application/pgp-keys}), import it into the isolated
    keyring.
  \item If no key was attached, export all public keys from the
    system keyring (\texttt{\~{}/.gnupg}) into the isolated keyring.
    This is a one-way bridge: read from system, never write.
  \item Run \texttt{gpg --verify} with the detached signature and
    signed data.
  \item Parse the exit code and stderr for validity, signer identity,
    and key ID.
  \item Remove the temporary directory.
\end{enumerate}

\subsection{Sign Flow}

\begin{enumerate}
  \item Check the signer key's expiry (\texttt{checkSignerKeyExpiry})
    before any other work --- keys with no expiration date, or an
    expiration date already passed, are rejected before boundary
    generation or body composition even begins.
  \item Compose the body as a MIME part with CRLF line endings
    (mandatory per RFC~3156).
  \item Write the GPG passphrase to a temporary file (mode 600) to
    avoid exposing it in process arguments or \texttt{ps} output.
  \item Run \texttt{gpg --detach-sign --armor} with
    \texttt{--passphrase-file}.
  \item Assemble the full RFC~822 message with a
    \texttt{multipart/signed} envelope containing the body part and
    detached signature.
  \item Generate the MIME boundary using \texttt{crypto/rand}.
\end{enumerate}

\subsection{Command-File Signature Verification}

Command-file verification (\texttt{internal/daemon/signature.go}'s
\texttt{VerifySignature}, DES-034/DES-035, both SETTLED) is a
structurally distinct, stricter path from \texttt{pgp.Verify} above:
it authorizes a \emph{command definition}, not an inbound message.

\begin{enumerate}
  \item Reconstruct the canonical bytes that were signed --- the
    command with its \texttt{Signature} field cleared, marshaled
    deterministically (\texttt{CanonicalCommandBytes}).
  \item Create an isolated GNUPGHOME, as with inbound verification.
  \item Import \emph{only} the authorizer's key (identified by a full
    40-hex fingerprint\footnote{\texttt{daemon.json}'s
    \texttt{owner\_handle}/\texttt{owner\_gpg\_key\_id} name this key
    as ``owner'' in code and config today; a rename to
    ``authorizer'' is proposed under bead beadle-4ul but not yet
    done.  This document uses ``authorizer'' for the concept and the
    current field names for the literal config keys.}) from the
    ambient keyring into the isolated one, and assert exactly one key
    with that exact fingerprint landed there --- never a silent
    ``verify against whatever's there.''
  \item Check that key's expiry against the isolated homedir (Design
    Invariant~4).
  \item Run \texttt{gpg --verify --status-fd} and parse the
    machine-readable \texttt{[GNUPG:]} status lines: every outcome
    except \texttt{GOODSIG} is an error, classified by reason
    (missing, invalid, wrong-key, key-expired).  This is a stricter
    parse than \texttt{pgp.Verify}'s, which reads a stderr substring.
\end{enumerate}

A nil error means the authorizer, and only the authorizer, signed this
exact command definition with a key that has not expired.
\texttt{LoadCommands}/\texttt{loadCommand} call \texttt{VerifySignature}
on every command file at daemon startup; when \texttt{daemon.json} is
absent, unresolvable, or names both or neither owner field, command
loading is disabled entirely --- the daemon starts with zero commands
rather than falling back to unsigned execution.

%%----------------------------------------------------------------------
\section{MCP Tool Surface}
%%----------------------------------------------------------------------

The tools below are registered by \texttt{internal/mcp/tools.go}'s
\texttt{RegisterTools} (message, contact, and identity tools; the last
21 rows unconditionally) plus \texttt{internal/mcp/poll\_tools.go}
(the final two rows, gated on a background poller being configured).
Re-run that lookup rather than trusting a stale copy of this table if
the tool surface has since changed --- the table is a snapshot, the
code is the source.

\begin{center}
\begin{longtable}{@{}llp{6cm}@{}}
\toprule
\textbf{Tool} & \textbf{Required Params} & \textbf{Description} \\
\midrule
\texttt{list\_messages}       & ---                            & List messages from a folder (id, from, date, subject, trust). Repo-scoped by default; \texttt{all\_repos} widens. \\
\texttt{search\_messages}     & ---                            & Search a folder by sender, subject, date, or free text; at least one filter required. Renders the same table as \texttt{list\_messages}. \\
\texttt{read\_message}        & \texttt{message\_id}           & Read a single message by UID: body, headers, attachment summary, trust level. \\
\texttt{list\_folders}        & ---                            & List all IMAP mailbox folders. \\
\texttt{send\_email}          & \texttt{to}, \texttt{subject}, \texttt{body} & Send via SMTP (primary) or Resend (fallback). Recipients may be contact names or aliases. \\
\texttt{reply\_message}       & \texttt{message\_id}, \texttt{body} & Reply to a message by UID, threaded (In-Reply-To/References), quoting the original. \\
\texttt{verify\_signature}    & \texttt{message\_id}           & Explicit PGP verification. Returns trust level, validity, key ID, signer. \\
\texttt{show\_mime}           & \texttt{message\_id}           & Show the MIME structure of a message. \\
\texttt{check\_trust}         & \texttt{message\_id}           & Classify trust level with a detailed explanation. \\
\texttt{move\_message}        & \texttt{message\_id}           & Move a message to another folder (default: Archive). \\
\texttt{batch\_move\_messages} & \texttt{message\_ids}          & Move multiple messages to another folder in one call. \\
\texttt{mark\_message}        & \texttt{message\_id}           & Set or clear the \textbackslash{}Seen flag; reading never changes read state. \\
\texttt{batch\_mark\_messages} & \texttt{message\_ids}          & Mark multiple messages read/unread in one call. \\
\texttt{download\_attachment} & \texttt{message\_id}, \texttt{part\_index} & Extract an attachment by MIME part index (from \texttt{show\_mime}). \\
\texttt{list\_contacts}       & ---                            & List all contacts in the address book. \\
\texttt{find\_contact}        & \texttt{query}                 & Look up a contact by name, email, or alias. \\
\texttt{add\_contact}         & \texttt{name}, \texttt{email}  & Add a contact; email may be a glob pattern, restricted to read-only or blocked permissions. \\
\texttt{remove\_contact}      & \texttt{name}                  & Remove a contact by name. \\
\texttt{whoami}               & ---                            & Show the active identity: email, handle, source, session participants. \\
\texttt{switch\_identity}     & ---                            & Switch the active identity for this session, or reset to default. \\
\texttt{enable}               & \texttt{action}                & Enable or disable beadle guidance (\texttt{CLAUDE.md} import, marker) in the current repo. \\
\texttt{set\_poll\_interval} (gated) & \texttt{interval}         & Set the background inbox polling interval. \\
\texttt{get\_poll\_status} (gated)   & ---                       & Show poller state: interval, active, last check time, unread count. \\
\bottomrule
\end{longtable}
\end{center}

Both gated rows are registered only when a background poller is
configured (\texttt{h.poller != nil}); the remaining rows above are
unconditional.

All tool handlers use a \texttt{withClient} pattern: a fresh IMAP
connection is established per request and closed when the handler
returns.  This avoids stale connection state at the cost of one
TCP handshake per tool call.

%%----------------------------------------------------------------------
\section{Design Invariants}
%%----------------------------------------------------------------------

These invariants are enforced by the architecture and hold across
all operations:

\begin{enumerate}
  \item \textbf{Zero agent authority (command execution only).}
    Every command the daemon executes requires a GPG-signed command
    file authorized by the owner (\S~Command-File Signature
    Verification) --- the daemon has no independent decision-making
    about \emph{what to run}.  This does not cover every action taken
    on inbound mail: \texttt{OnNewMail} creates a pipeline or a bare
    mission record from sender permission and transport trust alone,
    with no signed command anywhere in that earlier decision
    (\S~Daemon Architecture, Detection and Trust Gate).  With zero
    commands loaded, nothing downstream runs, but the mission/pipeline
    record itself is still created.

  \item \textbf{Preflight before execute.}  All permissions are
    validated before any command runs.  No partial execution.

  \item \textbf{Isolated verification.}  Inbound signature verification
    runs in a temporary GNUPGHOME so an untrusted sender's key never
    touches the system keyring. Signing and decryption use the owner's
    own key from the system keyring, since that is where it already
    lives.

  \item \textbf{Expired and non-expiring keys rejected.}  All
    command-signing keys must carry an expiration date that has not
    yet passed --- enforced on the primary key and, when the key has
    signing-capable subkeys, requiring at least one such subkey with a
    current, non-expired expiry (\texttt{internal/pgp/expiry.go}).

  \item \textbf{Audit log is tamperproof (designed, not
    implemented).}  Append-only, GPG-signed entries; only the owner
    can clear the log.  No code implements this today ---
    \texttt{internal/daemon} has no audit-log type, writer, or clear
    operation (bead beadle-623).  This invariant is retained as the
    design target, not a description of current behavior.

  \item \textbf{No secrets in logs.}  GPG key material, passwords,
    API keys, and raw email content are never logged.

  \item \textbf{No shell execution.}  All \texttt{exec.Command}
    calls use argument lists, never \texttt{shell=true}.

  \item \textbf{Header injection prevention.}  \texttt{ComposeRaw}
    and \texttt{Sign} reject header fields containing CR/LF.

  \item \textbf{File permission enforcement.}  Secret files must
    not be group or world readable (mode 600).

  \item \textbf{TLS verification skip is configurable, not
    loopback-only.}  Certificate verification is skipped
    automatically for loopback addresses (Proton Bridge's self-signed
    cert on localhost), but \texttt{tls\_skip\_verify} in config also
    disables it for \emph{any} host --- added for Docker host
    networking, where the bridge is not reachable via loopback.  The
    exception is a configuration choice, not a boundary enforced
    regardless of configuration.
\end{enumerate}

%%----------------------------------------------------------------------
\section{Security Boundaries}
%%----------------------------------------------------------------------

\subsection{Network Boundary}

\begin{itemize}
  \item \textbf{IMAP/SMTP}: \texttt{127.0.0.1} is the default host,
    not an enforced boundary --- \texttt{imap\_host}/\texttt{smtp\_host}
    are plain config fields, and \texttt{docs/TESTING.md}'s own
    Fastmail signing procedure switches \texttt{smtp\_host} to a
    non-loopback host.  STARTTLS (or implicit TLS on port 465,
    \S~Proton Bridge SMTP) with a self-signed-cert exception for
    loopback addresses; \texttt{tls\_skip\_verify} extends that
    exception to any host (Design Invariant~10).
  \item \textbf{Resend API}: outbound HTTPS to
    \texttt{api.resend.com}.  Bearer token authentication.
  \item \textbf{MCP}: stdio by default.  \texttt{serve --transport ws}
    instead binds \texttt{0.0.0.0} (all interfaces, not loopback) on
    port 8420 by default, with \texttt{CheckOrigin} accepting every
    origin unconditionally --- the only access control referenced
    anywhere in the codebase is an external \texttt{mcp-proxy}'s
    \texttt{MCP\_PROXY\_TOKEN}, named in a code comment, not enforced
    by this binary itself (bead beadle-nimy).
\end{itemize}

\subsection{Process Boundary}

\begin{itemize}
  \item \textbf{GPG}: subprocess. Verification uses an isolated
    GNUPGHOME; signing and decryption use the system keyring.
    Passphrase passed via temp file (mode 600), never via
    arguments or environment.
  \item \textbf{Keychain}: subprocess (\texttt{security} on macOS,
    \texttt{secret-tool} on Linux).  No direct library binding.
\end{itemize}

\subsection{Trust Boundary}

\begin{itemize}
  \item \textbf{Trusted}: Proton-to-Proton messages.  E2E
    encryption verified by Proton infrastructure, not by beadle.
  \item \textbf{Verified}: external messages with valid PGP
    signatures.  Verification performed by beadle in isolated GPG.
  \item \textbf{Untrusted}: external messages with invalid PGP
    signatures.  The signature was present but did not verify.
  \item \textbf{Unverified}: external messages with no signature.
    No cryptographic assertion about the sender.
\end{itemize}

%%----------------------------------------------------------------------
\section{Daemon Architecture}
%%----------------------------------------------------------------------

\texttt{beadle-daemon} is a standalone Go process that monitors email
autonomously and executes multi-stage pipelines when authorised
contacts send instructions.

\subsection{Detection and Trust Gate}

The daemon reuses the same \texttt{Poller} as \texttt{beadle-email},
decoupled via a \texttt{NewMailFunc} callback.  When new mail is
detected, the handler (\texttt{OnNewMail}, DES-027, SETTLED):

\begin{enumerate}
  \item Fetches recent unread messages via IMAP.
  \item Looks up the sender in the contact store.
  \item Checks the \texttt{x} (execute) permission bit.
  \item Verifies transport trust --- Proton E2E (trusted) or PGP
    signature with key bound to contact (verified).  Unverified
    messages from \texttt{rwx} contacts are rejected.  Trust in the
    Proton-E2E branch rests on IMAP actually being Proton Bridge on
    localhost, which nothing here confirms yet (bead beadle-hv7i).
  \item On pass: creates a pipeline when a spawner and command
    templates are configured; otherwise falls back to a bare mission
    record (\texttt{h.missions.Create}) with no pipeline or signed
    command involved at all (Design Invariant~1).
\end{enumerate}

\subsection{Pipeline Orchestrator}

A pipeline decomposes an email instruction into a sequence of typed
commands (DES-030, DES-031, both SETTLED).  Each command is a
GPG-signed YAML definition (\S~Command-File Signature Verification)
with a typed argument schema, write-set, budget, timeout, and tool
list, run by one of two runners:

\begin{itemize}
  \item \textbf{claude} --- spawns an ephemeral \texttt{claude -p
    --bare} session per stage (\S~Worker Isolation).
  \item \textbf{cli} --- runs a whitelisted binary directly, or a
    compound sequence of binaries declared as \texttt{steps}, chained
    together via \texttt{io.Pipe} goroutines.  A binary's resolved,
    symlink-followed path (\texttt{filepath.EvalSymlinks}) must still
    land inside a whitelisted directory.  Output is capped at 1~MB
    (\texttt{io.LimitReader}); anything beyond the cap is discarded
    and the result is marked truncated, not silently grown.
\end{itemize}

A stage's data flows to the next through a top-level \texttt{pipe}
value, shaped by each command's declared \texttt{mode}:

\begin{itemize}
  \item \texttt{process} --- the stage's output replaces \texttt{pipe}
    for the next stage, after validation against the command's
    \texttt{output\_schema} (\texttt{"text"} or a JSON Schema object).
  \item \texttt{passthrough} --- \texttt{pipe} is unchanged; the
    stage's own result is still recorded like any other stage.
\end{itemize}

Stage args and the prior stage's pipeline output are carried to the
worker as free text in the mission contract's top-level
\texttt{context} field --- ethos's \texttt{inputs} schema recognizes
only \texttt{trigger}, \texttt{files}, \texttt{ticket}, and
\texttt{references}; there is no \texttt{inputs.args} or
\texttt{inputs.previous\_output} field under any name.  Assuming
otherwise is exactly what shipped broken as beadle-8gt.

\begin{verbatim}
  Email instruction
    -> Planner (LLM or rule-based)
    -> [cmd1 (process/passthrough) | cmd2 | ... | auto-reply]
    -> Sequential mission execution; context field carries stage
       args + prior pipe value
\end{verbatim}

File-flow support --- attachments and generated artifacts moving
between stages as file references, rather than only through
\texttt{pipe} --- is designed (DES-032) but not implemented:
\texttt{Command} has no file-flow fields today.

\subsection{Auto-Reply (DES-029)}

The executor auto-appends a \texttt{reply} command as the terminal
stage of every pipeline.  The planner returns work stages only; the
reply is the daemon's I/O framing.  This keeps commands
channel-agnostic --- different channels (email, chat) use different
reply commands without modifying the pipeline logic.

On failure, the else handler sends a fixed-text error reply.  The
originator never receives silence.

\subsection{Worker Isolation}

Isolation here is process-level only, not container-level: workers
share the daemon host's kernel, network namespace, and filesystem
outside their restricted \texttt{PATH}.  Container isolation (network
namespace, seccomp) is an unfiled TODO in \texttt{spawner.go}, not yet
implemented.  Within that scope, workers run with:

\begin{itemize}
  \item \texttt{--bare} mode (no hooks, plugins, CLAUDE.md).
  \item Restricted PATH: the resolved \texttt{claude} binary's own
    directory, plus \texttt{/usr/bin} and \texttt{/usr/local/bin}.
  \item Per-command MCP servers from a registry.
  \item Adversarial robustness instructions in the system prompt.
  \item 30-minute timeout, \$5 budget cap, 50-turn limit --- each
    overridable; these are the defaults.
  \item Protected env vars (HOME, PATH, ANTHROPIC\_API\_KEY) forced
    after command overrides.
\end{itemize}

\subsection{Pipeline Persistence}

Pipeline state is written to
\texttt{\~{}/.punt-labs/beadle/pipelines/<id>.json} via atomic rename
(write to \texttt{.tmp-<id>.json}, then \texttt{os.Rename}) at each
state transition, under 0700/0600 permissions, with a path-traversal
guard on the pipeline ID.  On daemon startup, in-flight
(\texttt{"running"}) pipelines are marked failed and the originator is
notified.  A 30-day retention sweep runs on a recurring ticker
(\texttt{DefaultRetention}); orphaned \texttt{.tmp-*.json} files left
by a process that died between write and rename are reaped separately,
after 5 minutes, since a torn write may not even parse as a pipeline
record.  \texttt{cmd/beadle-daemon}'s \texttt{main} wires this store
into both the startup sweep and the \texttt{Executor} that runs live
pipelines, so both see the same records.

%%----------------------------------------------------------------------
\section{Key Properties}
%%----------------------------------------------------------------------

The specification captures these properties, verifiable by
inspection and testing:

\begin{enumerate}
  \item \textbf{Trust totality}: every inbound message receives
    exactly one of the four trust levels.  The
    \texttt{ClassifyTrust} schema is total over all input
    combinations.

  \item \textbf{Credential priority}: keychain always takes
    priority over file, file over environment variable.  The
    \texttt{ResolveCredential} schema enforces this ordering.

  \item \textbf{GPG isolation (verification only)}: the $isolated =
    ztrue$ invariant in \texttt{GPGEnvironment} ensures no PGP
    \emph{verification} operation writes to the user's system
    keyring.  Signing and decryption are a deliberate exception ---
    they use the owner's own key from the system keyring, since that
    is where it already lives (Design Invariant~3, \S~PGP Signing,
    DES-024).

  \item \textbf{Send fallback}: if SMTP is available, it is
    always preferred.  Resend is used only when SMTP is
    unreachable.

  \item \textbf{Connection-per-request}: each MCP tool call
    creates and destroys its own IMAP connection, preventing
    stale state across calls.

  \item \textbf{Trust gate before creation}: only messages with
    transport trust $\geq verified$ (PGP) or $= trusted$ (Proton E2E)
    AND contact \texttt{x}-bit reach mission or pipeline creation
    (\S~Detection and Trust Gate).

  \item \textbf{Auto-reply}: every pipeline terminates with a reply
    to the originator.  Success sends the final output.  Failure
    sends a fixed-text error with a pipeline reference ID.

  \item \textbf{Stage isolation}: each pipeline stage runs as a
    separate \texttt{claude -p --bare} process with per-command
    tools, budget, and timeout.  Missions are closed between stages
    to release write-set locks.

  \item \textbf{Crash recovery}: pipeline state is persisted at each
    transition.  On restart, in-flight pipelines are marked failed
    and originators are notified.
\end{enumerate}

\end{document}
